\documentclass[11pt]{article}

\usepackage[utf8]{inputenc}
\usepackage[T1]{fontenc}
\usepackage[spanish,es-noshorthands]{babel}
\usepackage[margin=2.5cm]{geometry}
\usepackage{graphicx}
\usepackage{booktabs}
\usepackage{longtable}
\usepackage{amsmath}
\usepackage{listings}
\usepackage{xcolor}
\usepackage{hyperref}
\usepackage{tikz}
\usetikzlibrary{arrows.meta,positioning,shapes.geometric,calc}
\usepackage{caption}

\hypersetup{colorlinks=true,linkcolor=black,urlcolor=blue,citecolor=black}

% Estilo para listados de codigo C.
\definecolor{codegray}{gray}{0.95}
\definecolor{kw}{RGB}{0,0,160}
\definecolor{cmt}{RGB}{0,120,0}
\definecolor{str}{RGB}{160,0,0}
\lstset{
  language=C, basicstyle=\ttfamily\footnotesize, keywordstyle=\color{kw},
  commentstyle=\color{cmt}\itshape, stringstyle=\color{str},
  backgroundcolor=\color{codegray}, frame=single, breaklines=true,
  showstringspaces=false, numbers=left, numberstyle=\tiny\color{gray},
  tabsize=2, captionpos=b
}

% Incluye una figura solo si el archivo existe (permite compilar sin las graficas aun).
\newcommand{\figsi}[3]{%
  \IfFileExists{#1}{%
    \begin{figure}[ht]\centering
      \includegraphics[width=#2\linewidth]{#1}
      \caption{#3}\end{figure}}%
    {\begin{center}\fbox{\parbox{0.8\linewidth}{\centering\small\textit{[Figura pendiente: #1. Ejecute \texttt{scripts/run\_bench.sh} y \texttt{scripts/plot.py}.]}\\#3}}\end{center}}}

\begin{document}

% ====================== CARATULA ======================
\begin{titlepage}
\centering
\vspace*{1.5cm}
{\large Universidad del Valle de Guatemala\par}
{\large Facultad de Ingeniería\par}
{\large Departamento de Ciencias de la Computación\par}
{\large CC3069 Computación Paralela y Distribuida -- Sección 20, Semestre II 2026\par}
\vspace{2.5cm}
{\Huge\bfseries Proyecto \#2\par}
\vspace{0.6cm}
{\LARGE Fuerza bruta de llaves DES con Open MPI\par}
\vspace{2.5cm}
{\large\bfseries Integrante:\par}
\vspace{0.3cm}
{\large Luis Palacios \#23933\par}
\vfill
{\large \today\par}
\end{titlepage}

\tableofcontents
\newpage

% ====================== INTRODUCCION ======================
\section{Introducción}
Este proyecto aplica el modelo de memoria distribuida con \textbf{Open MPI}~\cite{mpi40,pacheco2011}
al problema de romper, por \emph{fuerza bruta}, la llave privada con la que se cifró un texto usando
el algoritmo \textbf{DES}~\cite{fips463,stallings2017}. El criterio para reconocer una llave correcta
es que el texto descifrado contenga, como subcadena, una \emph{frase clave} conocida de antemano
(por ejemplo \texttt{"es una prueba de"}). Se eligió una frase larga para que sea muy improbable que
aparezca por azar al descifrar con una llave equivocada.

Además de implementar el cifrado/descifrado y la paralelización, el objetivo central es
\textbf{analizar el speedup} del algoritmo y, sobre todo, entender por qué los tiempos paralelos del
enfoque ingenuo (\emph{naive}) son tan inconsistentes y dependientes de la llave. A partir de ese
análisis se proponen e implementan tres acercamientos alternativos y una versión híbrida MPI+OpenMP.

El código es de autoría propia en C. Como la cabecera \texttt{rpc/des\_crypt.h} del código base fue
removida de las versiones modernas de \texttt{glibc}, se implementó un \textbf{DES propio y
auto-contenido} (\texttt{src/des.c}) que reproduce la misma API (\texttt{ecb\_crypt},
\texttt{des\_setparity}) para no alterar la lógica del resto del programa.

% ====================== PARTE A ======================
\section{DES: fundamentos y pasos de cifrado/descifrado}
El \emph{Data Encryption Standard} (DES)~\cite{fips463} es un cifrador por bloques de clave simétrica
que opera sobre bloques de \textbf{64 bits} con una llave de 56 bits efectivos (8 bytes, de los cuales
un bit por byte es de paridad). Es una red de Feistel de 16 rondas. Los pasos para cifrar un bloque son:

\begin{enumerate}
  \item \textbf{Permutación inicial (IP):} se reordenan los 64 bits del bloque de entrada.
  \item \textbf{División:} el bloque permutado se parte en dos mitades de 32 bits, $L_0$ y $R_0$.
  \item \textbf{16 rondas de Feistel:} para $i = 1,\dots,16$:
    \[ L_i = R_{i-1}, \qquad R_i = L_{i-1} \oplus f(R_{i-1}, K_i), \]
    donde $K_i$ es la subllave de 48 bits de la ronda $i$ y $f$ es la función de ronda.
  \item \textbf{Intercambio e IP$^{-1}$:} se concatena $R_{16}L_{16}$ y se aplica la permutación final
    (inversa de IP) para obtener el texto cifrado.
\end{enumerate}

La \textbf{función de ronda} $f(R,K)$ realiza: (a) \emph{expansión} de $R$ de 32 a 48 bits (tabla E);
(b) XOR con la subllave $K_i$; (c) sustitución mediante las 8 \textbf{S-boxes} (cada grupo de 6 bits
produce 4 bits, usando los bits externos como fila y los internos como columna); y (d) una
\emph{permutación} P de 32 bits.

El \textbf{calendario de llaves} aplica PC-1 a la llave de 64 bits (descartando los 8 bits de paridad),
divide el resultado de 56 bits en $C_0$ y $D_0$ de 28 bits, y en cada ronda los rota a la izquierda
(1 o 2 posiciones según la ronda) antes de aplicar PC-2 para obtener la subllave $K_i$ de 48 bits.

El \textbf{descifrado} es idéntico al cifrado pero aplicando las subllaves en orden inverso
($K_{16},\dots,K_1$); esta es la elegante propiedad de las redes de Feistel.

\subsection{Diagrama de flujo de DES}
\begin{center}
\begin{tikzpicture}[node distance=7mm and 10mm,
  box/.style={draw,rounded corners,align=center,minimum width=3.2cm,minimum height=7mm,font=\small},
  op/.style={draw,circle,inner sep=1pt,font=\small},
  arr/.style={-{Latex[length=2mm]}}]
  \node[box] (in) {Bloque de 64 bits};
  \node[box,below=of in] (ip) {Permutación inicial (IP)};
  \node[box,below=of ip] (split) {$L_0 \;\|\; R_0$ (32 + 32 bits)};
  \node[box,below=of split,fill=codegray] (round) {16 rondas de Feistel\\ $L_i=R_{i-1}$, \; $R_i=L_{i-1}\oplus f(R_{i-1},K_i)$};
  \node[box,below=of round] (swap) {Intercambio $R_{16}\|L_{16}$};
  \node[box,below=of swap] (fp) {Permutación final (IP$^{-1}$)};
  \node[box,below=of fp] (out) {Texto cifrado (64 bits)};
  \draw[arr] (in)--(ip); \draw[arr] (ip)--(split); \draw[arr] (split)--(round);
  \draw[arr] (round)--(swap); \draw[arr] (swap)--(fp); \draw[arr] (fp)--(out);

  % Detalle de la funcion f a la derecha.
  \node[box,right=2.4cm of round,yshift=1.3cm] (e) {Expansión E (32$\to$48)};
  \node[op,below=5mm of e] (xor) {$\oplus$};
  \node[box,right=6mm of xor,font=\footnotesize] (k) {$K_i$ (48)};
  \node[box,below=5mm of xor] (s) {8 S-boxes (48$\to$32)};
  \node[box,below=5mm of s] (p) {Permutación P (32)};
  \draw[arr] (e)--(xor); \draw[arr] (k)--(xor); \draw[arr] (xor)--(s); \draw[arr] (s)--(p);
  \node[font=\small] at ($(e)+(0,0.7)$) {\textbf{Función } $f(R,K_i)$};
\end{tikzpicture}
\end{center}

\section{Rutinas del programa base}
Se explican las rutinas del \texttt{bruteforce.c} base (y sus equivalentes en nuestra implementación).

\paragraph{\texttt{encrypt(key, *ciph, len)} y \texttt{decrypt(key, *ciph, len)}.}
Ambas reciben una llave \texttt{long}, la \emph{expanden} a los 8 bytes que DES necesita y llaman a
\texttt{ecb\_crypt} en modo cifrar/descifrar. La expansión toma 7 bits de la llave por cada byte y deja
el bit menos significativo para la paridad (que \texttt{des\_setparity} ajusta a paridad impar). El
código base hacía esta expansión con \texttt{0xFE << i*8} usando \texttt{int}, lo que \textbf{desborda}
para los bytes altos ($i\ge 4$); por eso su comentario advertía que la llave debía ser pequeña. En
nuestra versión (\texttt{spread\_key}) se corrige usando corrimientos sobre \texttt{long}, válido para
todo el rango $[0,2^{56})$:

\begin{lstlisting}[caption={Expansión de llave corregida (src/common.c).}]
void spread_key(long key, unsigned char out8[8]) {
  for (int i = 0; i < 8; ++i) {
    unsigned char seven = (unsigned char)((key >> (7 * i)) & 0x7F);
    out8[i] = (unsigned char)(seven << 1);  // 7 bits + LSB para paridad
  }
  des_setparity((char *)out8);
}
\end{lstlisting}

\paragraph{\texttt{tryKey(key, *ciph, len)}.}
Copia el texto cifrado, lo descifra con la llave candidata y verifica, con \texttt{strstr}, si contiene
la frase clave. Devuelve 1 (verdadero) si la llave parece correcta. Es el núcleo del lazo de búsqueda.

\paragraph{\texttt{memcpy(dest, src, n)}.}
Copia \texttt{n} bytes de \texttt{src} a \texttt{dest}. Se usa para trabajar sobre una copia del texto
cifrado sin destruir el original (el descifrado es \emph{in situ}).

\paragraph{\texttt{strstr(haystack, needle)}.}
Busca la primera aparición de la subcadena \texttt{needle} dentro de \texttt{haystack}; devuelve un
puntero a ella o \texttt{NULL}. Es el criterio de aceptación de una llave.

\section{Primitivas MPI y flujo de comunicación}
La búsqueda paralela usa las siguientes primitivas~\cite{mpi40,pacheco2011}:

\paragraph{\texttt{MPI\_Irecv}.} Recepción \textbf{no bloqueante}: publica una solicitud para recibir el
anuncio ``llave encontrada'' desde cualquier proceso, sin detener la búsqueda. Devuelve un
\texttt{MPI\_Request} que luego se consulta.

\paragraph{\texttt{MPI\_Send}.} Envío (bloqueante) del valor de la llave encontrada. Cuando un proceso
halla la llave, hace \texttt{MPI\_Send} a \emph{todos} los procesos para avisarles que se detengan.

\paragraph{\texttt{MPI\_Wait}.} Espera a que una solicitud no bloqueante (la de \texttt{MPI\_Irecv})
se complete. Se usa al final para consumir limpiamente el mensaje pendiente.

\paragraph{Correcciones al flujo del base.} El \texttt{bruteforce.c} original publicaba el
\texttt{MPI\_Irecv} pero \textbf{nunca} comprobaba su estado dentro del lazo, de modo que los procesos
que no encontraban la llave jamás se enteraban del hallazgo ajeno y recorrían todo su rango. Nuestra
versión añade \texttt{MPI\_Test} periódico para que \emph{todos} se detengan al encontrarse la llave, y
usa \texttt{MPI\_Bcast} (difusión del texto cifrado desde el rank 0) y \texttt{MPI\_Allreduce} (acuerdo
final sobre la llave global) para un cierre robusto.

El flujo es: el rank 0 lee y valida el archivo y lo \emph{difunde} (\texttt{MPI\_Bcast}); cada proceso
publica un \texttt{MPI\_Irecv} y recorre su porción del espacio; cada cierto número de iteraciones hace
\texttt{MPI\_Test}; el que encuentra la llave hace \texttt{MPI\_Send} a todos; finalmente un
\texttt{MPI\_Allreduce} consolida el resultado y el rank 0 lo imprime.

\section{Versión secuencial (Parte A)}
La versión secuencial (\texttt{src/secuencial.c}) recorre el espacio $[0,\text{cota})$ probando cada
llave con \texttt{try\_key} hasta hallar la frase clave, midiendo el tiempo con un reloj monótono. Sirve
como línea base para el \textbf{speedup} ($S = T_\text{secuencial}/T_\text{paralelo}$) y para observar
cómo crece el tiempo con la magnitud de la llave: el costo es proporcional a la \emph{posición} de la
llave en el recorrido, pues la búsqueda se detiene al encontrarla.

\section{Cifrado desde archivo (Parte B1)}
La herramienta \texttt{src/cifrar.c} carga un texto desde un archivo \texttt{.txt}, lo rellena con ceros
hasta un múltiplo de 8 bytes (tamaño de bloque DES) y lo cifra con una llave arbitraria pasada como
parámetro; también puede descifrar. La Figura/registro de Anexo~2 evidencia el cifrado y descifrado de
una oración sencilla con la llave \texttt{42}, recuperando exactamente el texto original.

\begin{lstlisting}[language=bash,caption={Demostración de cifrado/descifrado con llave 42.}]
$ printf "Esta es una prueba de proyecto 2" > entrada.txt
$ ./bin/cifrar cifrar entrada.txt cifrado.bin 42
Cifrado OK: 'entrada.txt' (32 bytes) -> 'cifrado.bin' (32 bytes) con llave 42.
$ ./bin/cifrar descifrar cifrado.bin salida.txt 42
Descifrado OK: 'cifrado.bin' -> 'salida.txt' con llave 42.
Texto recuperado: "Esta es una prueba de proyecto 2"
\end{lstlisting}

\section{Pruebas de la Parte B2 y reflexión}
Con 4 procesos, el texto \texttt{"Esta es una prueba de proyecto 2"} y la frase \texttt{"es una prueba
de"}, se midieron las llaves solicitadas:
\begin{itemize}
  \item \textbf{(a) llave \texttt{123456}:} se encuentra muy rápido (posición baja del espacio).
  \item \textbf{(b) llave $2^{56}/4 = 18014398509481983$:} está al \emph{final} del primer bloque en
    el reparto naive de 4 procesos; equivale a recorrer $\approx 2^{54}$ llaves, por lo que
    \textbf{no termina} en tiempo razonable (como anticipa el enunciado).
  \item \textbf{(c) llave $2^{56}/4 + 1 = 18014398509481984$:} está al \emph{inicio} del segundo
    bloque; el proceso \#2 la encuentra en su primera iteración, dando un speedup aparente enorme.
\end{itemize}

\paragraph{Reflexión (B2.d).} El tiempo paralelo del naive \textbf{no depende solo de la dificultad de
la llave sino de su posición dentro del reparto por bloques}. Dos llaves numéricamente contiguas
producen tiempos opuestos: con \texttt{-np 4} la llave $2^{56}/4-1=18014398509481983$ (final del bloque
0) \textbf{no termina}, mientras que $2^{56}/4=18014398509481984$ (inicio del bloque 1) se encuentra en
$\approx 0.05$~s. Esto muestra que el speedup del naive es \emph{engañoso}: la rapidez de la segunda no
viene del paralelismo sino de la \emph{casualidad} de que la llave quedó justo al inicio de un bloque.
La llave \texttt{123456}, en cambio, cae en el bloque 0 y el naive la busca igual que el secuencial (sin
speedup), mientras que el enfoque cíclico sí reparte ese trabajo entre los procesos. Todo esto motiva
repartos que no dependan de dónde cae la llave.

\section{Acercamientos propuestos (Partes B3--B4)}
Sea $K$ la posición de la llave y $N$ el número de procesos.

\paragraph{Naive (base).} Divide $[0,\text{cota})$ en $N$ bloques contiguos. El tiempo paralelo depende
del bloque donde cae $K$ y de su posición dentro de él: el mejor caso es $O(1)$ (inicio de un bloque) y
el peor es $O(\text{cota}/N)$ con speedup $\approx 1$ (final del bloque 0). \emph{Tiempo esperado muy
variable.}

\paragraph{Enfoque 1 -- Cíclico (\emph{stride}).} El proceso \texttt{id} prueba las llaves
$id, id{+}N, id{+}2N,\dots$ El espacio avanza de forma globalmente uniforme, así que la llave en la
posición $K$ se alcanza tras $\approx K/N$ iteraciones \textbf{sin importar su valor}. El tiempo
paralelo esperado es $\approx K/N$ de manera \emph{consistente}, eliminando el peor caso del naive.

\paragraph{Enfoque 2 -- Maestro--trabajador (reparto dinámico).} El rank 0 reparte ``chunks'' del
espacio bajo demanda; cada trabajador pide, procesa y vuelve a pedir. Aporta \textbf{balanceo de carga}
y robustez cuando los procesos avanzan a ritmos distintos (máquinas heterogéneas), a cambio de dedicar
un proceso a coordinar y de cierto costo de comunicación (controlable con el tamaño de chunk).

\paragraph{Enfoque 3 -- Aleatorio (biyección Feistel).} Cada proceso recorre su parte del espacio en
orden \emph{pseudoaleatorio} usando una biyección (red de Feistel) que garantiza cubrir cada llave
exactamente una vez. El tiempo esperado se vuelve \textbf{insensible al valor} de la llave: ya no hay
llaves sistemáticamente ``fáciles'' o ``difíciles'' según su magnitud, solo según el azar del orden.

\section{Resultados: speedup y consistencia (Parte B5)}
Las pruebas se realizaron en WSL2 Ubuntu 24.04 (GCC 13.3, Open MPI) sobre un espacio acotado y medible
($\text{cota}=2^{23}$) para obtener speedups reales; las definiciones de llave fácil/media/difícil del
enunciado (fracciones de $2^{56}$) se replican como fracciones de la cota. Cada punto es la mediana de
varias repeticiones (ver \texttt{scripts/run\_bench.sh}).

\figsi{img/naive_inconsistencia.png}{0.75}{Inconsistencia del naive según la posición de la llave:
mejor caso (inicio de bloque) vs. peor caso (final del bloque 0), con 4 procesos. El cíclico se mantiene
estable en ambos.}

\figsi{img/speedup_media.png}{0.8}{Speedup vs. número de procesos (escenario ``media''). El cíclico y el
aleatorio escalan de forma consistente; el naive depende del caso.}

\figsi{img/speedup_dificil.png}{0.8}{Speedup vs. número de procesos (escenario ``difícil'').}

\figsi{img/tiempo_media.png}{0.8}{Tiempo de búsqueda vs. número de procesos (escenario ``media'').}

La Tabla~\ref{tab:medianas} resume las medianas observadas (en segundos) con \texttt{-np 4} en nuestro
equipo. Se aprecia el fenómeno central:

\begin{table}[ht]\centering
\caption{Mediana de tiempo (s) por escenario y programa, \texttt{-np 4}, cota $=2^{20}$.
``T.O.'' = no terminó dentro del límite.}
\label{tab:medianas}
\begin{tabular}{lrrrrr}
\toprule
Escenario & $T_\text{seq}$ & naive & cíclico & maestro--trab. & aleatorio \\
\midrule
naive\_mejor (inicio de bloque) & 2.45 & \textbf{0.04} & 0.90 & 1.50 & 1.27 \\
naive\_peor (final del bloque 0) & 2.42 & \textbf{3.00} & 0.80 & 1.58 & 2.08 \\
fácil   & 4.91 & 0.04 & 1.60 & 2.61 & 1.64 \\
media   & 6.04 & 1.54 & 2.01 & 2.76 & 1.12 \\
difícil & 2.26 & 3.30 & 0.73 & 1.66 & \textbf{0.43} \\
\bottomrule
\end{tabular}
\end{table}

\begin{itemize}
  \item \textbf{El naive es errático e incluso contraproducente.} En \texttt{naive\_mejor} es casi
    instantáneo (0.04~s, speedup aparente $\sim$60$\times$) solo porque la llave cae al inicio de un
    bloque; pero en \texttt{naive\_peor} tarda \emph{más} que el secuencial (3.00~s vs 2.42~s,
    \textbf{speedup $<1$}), porque los procesos 1--3 hacen trabajo inútil y le compiten CPU al proceso 0,
    que es el único que avanza hacia la llave.
  \item \textbf{El cíclico es consistente}: $\sim$0.7--2.0~s en todos los escenarios (speedup real
    $\approx 3\times$ con 4 procesos), sin mejores ni peores casos patológicos.
  \item \textbf{El aleatorio} es insensible al valor de la llave; brilla cuando la llave está ``lejos''
    numéricamente (en \texttt{difícil} fue el más rápido, 0.43~s, speedup $\approx 5\times$), aunque su
    tiempo tiene más varianza por ser un recorrido al azar.
  \item \textbf{El maestro--trabajador} escala de forma estable, con una penalización por dedicar el
    rank 0 a coordinar: con \texttt{-np 4} solo hay 3 trabajadores (y con \texttt{-np 2}, uno solo, por
    lo que no acelera). Su fortaleza real es el balanceo ante nodos heterogéneos.
\end{itemize}

Las mediciones se hicieron en un solo equipo con \emph{oversubscribe} (más procesos que núcleos
físicos), lo que explica que a veces agregar procesos no ayude o incluso perjudique cuando el trabajo no
está bien repartido (caso del naive).

% ====================== EXTRA ======================
\section{Extra: híbrido MPI + OpenMP}
La versión \texttt{src/hibrido.c} combina memoria distribuida (MPI entre procesos) y memoria compartida
(OpenMP entre hilos dentro de cada proceso)~\cite{barlas2015,openmp}. Cada proceso toma su carril
cíclico y lo reparte entre sus hilos; MPI se inicializa en modo \texttt{MPI\_THREAD\_FUNNELED} (solo el
hilo principal hace llamadas MPI). Con $T$ hilos por proceso el paralelismo efectivo es $N\cdot T$.

\figsi{img/speedup_facil.png}{0.8}{Escalamiento con el enfoque híbrido y los enfoques base (escenario
``fácil'').}

Se observa un speedup adicional al aumentar los hilos, consistente con la versión sin OpenMP. En el
escenario ``media'' (cota $=2^{20}$, $T_\text{seq}=6.04$~s), con 1 proceso el tiempo pasó de 6.12~s
(1 hilo) a 3.13~s (4 hilos), y con 4 procesos de 1.90~s (1 hilo) a 1.21~s (4 hilos): el paralelismo de
hilos se suma al de procesos. La ganancia marginal se reduce cuando no hay suficientes núcleos físicos,
pues el \emph{hyper-threading} y la contención de recursos limitan el escalamiento (como también se
documentó para este equipo en laboratorios previos del curso).

% ====================== DISCUSION ======================
\section{Discusión, conclusiones y recomendaciones}
\begin{itemize}
  \item El reto principal \textbf{no} fue paralelizar, sino entender por qué el speedup del naive es
    engañoso: depende de la \emph{posición} de la llave en el reparto, no solo de su dificultad.
  \item Repartir el espacio de forma cíclica o pseudoaleatoria produce tiempos paralelos
    \textbf{consistentes} ($\approx K/N$), que reflejan mejor el beneficio real del paralelismo.
  \item Se corrigieron dos defectos del código base: el desbordamiento de \texttt{int} en la expansión
    de llave y la terminación temprana rota (faltaba \texttt{MPI\_Test}).
  \item El maestro--trabajador es la mejor opción cuando los nodos son heterogéneos; el cíclico es el
    más simple y eficiente en nodos homogéneos; el aleatorio es el más robusto ante llaves adversarias.
  \item \textbf{Recomendaciones:} ajustar el tamaño de chunk según la latencia de red; usar el híbrido
    MPI+OpenMP para aprovechar todos los núcleos; y medir siempre con varias llaves para no sacar
    conclusiones de un speedup aparente.
\end{itemize}

% ====================== ANEXO 1 ======================
\newpage
\appendix
\section{Anexo 1 -- Catálogo de funciones y librerías}

\paragraph{\texttt{des.c / des.h} (DES propio).}
\begin{itemize}
  \item \texttt{int ecb\_crypt(char *key, char *data, unsigned datalen, unsigned mode)} --
    \emph{Entradas:} llave de 8 bytes, buffer de datos, longitud (múltiplo de 8), modo
    (\texttt{DES\_ENCRYPT}/\texttt{DES\_DECRYPT}). \emph{Salida:} \texttt{data} cifrado/descifrado
    \emph{in situ}; retorna 0 en éxito. \emph{Propósito:} cifrado/descifrado ECB bloque a bloque.
  \item \texttt{void des\_setparity(char *key)} -- \emph{Entrada/Salida:} llave de 8 bytes; ajusta el
    LSB de cada byte a paridad impar.
  \item \texttt{uint64\_t des\_crypt\_block(uint64\_t block, uint64\_t key64, int decrypt)} --
    \emph{Entradas:} bloque y llave de 64 bits, bandera de descifrado. \emph{Salida:} bloque procesado.
    \emph{Propósito:} núcleo de 16 rondas de Feistel (expuesto para pruebas).
\end{itemize}

\paragraph{\texttt{common.c / common.h} (utilidades).}
\begin{itemize}
  \item \texttt{void spread\_key(long key, unsigned char out8[8])} -- expande una llave de 56 bits a 8
    bytes con paridad.
  \item \texttt{void decrypt\_key/encrypt\_key(long key, char *ciph, int len)} -- cifra/descifra con una
    llave \texttt{long}.
  \item \texttt{int try\_key(long key, const unsigned char *ciph, int len, const char *search)} --
    \emph{Salida:} 1 si al descifrar aparece la frase \texttt{search}. Núcleo de la búsqueda.
  \item \texttt{unsigned char *read\_file(const char *path, long *len)} / \texttt{int write\_file(...)}
    -- E/S binaria de archivos.
  \item \texttt{long pad\_len(long)} -- redondeo al múltiplo de 8. \quad
    \texttt{double now\_seconds(void)} -- reloj monótono. \quad
    \texttt{int parse\_long(const char*, long*)} -- parseo defensivo.
\end{itemize}

\paragraph{\texttt{mpi\_search.h} (andamiaje MPI).}
\begin{itemize}
  \item \texttt{load\_cipher\_bcast(path, *len, rank)} -- el rank 0 lee/valida y difunde el cifrado.
  \item \texttt{print\_result(...)} -- imprime llave, archivo, frase y texto descifrado.
\end{itemize}

\paragraph{Programas.} \texttt{cifrar}, \texttt{secuencial}, \texttt{naive}, \texttt{ciclico},
\texttt{maestro\_trabajador}, \texttt{aleatorio}, \texttt{hibrido}, \texttt{des\_test}.
\paragraph{Librerías.} Open MPI (\texttt{mpi.h}), OpenMP (\texttt{omp.h}), C estándar. No se usa
\texttt{rpc/des\_crypt.h} ni OpenSSL.

% ====================== ANEXO 2 ======================
\section{Anexo 2 -- Bitácora de pruebas y speedups}
Tabla de tiempos (mediana de repeticiones) por escenario, programa y número de procesos. Las entradas
\texttt{TIMEOUT} corresponden a llaves que no terminan en el tiempo límite (p.\,ej. $2^{56}/4$).

\IfFileExists{tabla_bitacora.tex}{% Tabla generada automaticamente por scripts/plot.py
\begin{longtable}{llrrr}
\toprule
Escenario & Programa & np & hilos & Tiempo (s) / estado \\
\midrule
\endhead
b2\_123456 & aleatorio & 4 & 1 & TIMEOUT \\
b2\_123456 & ciclico & 4 & 1 & 0.457455 \\
b2\_123456 & maestro\_trabajador & 4 & 1 & 0.723570 \\
b2\_123456 & naive & 4 & 1 & 1.444871 \\
b2\_123456 & secuencial & 1 & 1 & 1.194581 \\
b2\_18014398509481983 & aleatorio & 4 & 1 & TIMEOUT \\
b2\_18014398509481983 & ciclico & 4 & 1 & TIMEOUT \\
b2\_18014398509481983 & maestro\_trabajador & 4 & 1 & TIMEOUT \\
b2\_18014398509481983 & naive & 4 & 1 & TIMEOUT \\
b2\_18014398509481983 & secuencial & 1 & 1 & TIMEOUT \\
b2\_18014398509481984 & aleatorio & 4 & 1 & TIMEOUT \\
b2\_18014398509481984 & ciclico & 4 & 1 & TIMEOUT \\
b2\_18014398509481984 & maestro\_trabajador & 4 & 1 & TIMEOUT \\
b2\_18014398509481984 & naive & 4 & 1 & 0.047454 \\
b2\_18014398509481984 & secuencial & 1 & 1 & TIMEOUT \\
dificil & aleatorio & 1 & 1 & 1.573220 \\
dificil & aleatorio & 2 & 1 & 0.793282 \\
dificil & aleatorio & 4 & 1 & 0.432915 \\
dificil & ciclico & 1 & 1 & 2.215836 \\
dificil & ciclico & 2 & 1 & 1.269111 \\
dificil & ciclico & 4 & 1 & 0.728328 \\
dificil & maestro\_trabajador & 1 & 1 & 2.142168 \\
dificil & maestro\_trabajador & 2 & 1 & 2.296880 \\
dificil & maestro\_trabajador & 4 & 1 & 1.657954 \\
dificil & naive & 1 & 1 & 2.228230 \\
dificil & naive & 2 & 1 & 2.425331 \\
dificil & naive & 4 & 1 & 3.299544 \\
dificil & secuencial & 1 & 1 & 2.260918 \\
facil & aleatorio & 1 & 1 & 5.728338 \\
facil & aleatorio & 2 & 1 & 2.983694 \\
facil & aleatorio & 4 & 1 & 1.635906 \\
facil & ciclico & 1 & 1 & 4.899763 \\
facil & ciclico & 2 & 1 & 2.579342 \\
facil & ciclico & 4 & 1 & 1.595177 \\
facil & maestro\_trabajador & 1 & 1 & 4.908384 \\
facil & maestro\_trabajador & 2 & 1 & 5.249495 \\
facil & maestro\_trabajador & 4 & 1 & 2.610317 \\
facil & naive & 1 & 1 & 4.922657 \\
facil & naive & 2 & 1 & 0.039670 \\
facil & naive & 4 & 1 & 0.043740 \\
facil & secuencial & 1 & 1 & 4.911791 \\
media & aleatorio & 1 & 1 & 3.277386 \\
media & aleatorio & 2 & 1 & 1.739608 \\
media & aleatorio & 4 & 1 & 1.120650 \\
media & ciclico & 1 & 1 & 6.149324 \\
media & ciclico & 2 & 1 & 3.247409 \\
media & ciclico & 4 & 1 & 2.007713 \\
media & hibrido & 1 & 1 & 6.123998 \\
media & hibrido & 1 & 2 & 3.167114 \\
media & hibrido & 1 & 4 & 3.126788 \\
media & hibrido & 2 & 1 & 3.186899 \\
media & hibrido & 2 & 2 & 1.815153 \\
media & hibrido & 2 & 4 & 2.107325 \\
media & hibrido & 4 & 1 & 1.900100 \\
media & hibrido & 4 & 2 & 1.583883 \\
media & hibrido & 4 & 4 & 1.213953 \\
media & maestro\_trabajador & 1 & 1 & 6.104929 \\
media & maestro\_trabajador & 2 & 1 & 6.181267 \\
media & maestro\_trabajador & 4 & 1 & 2.763159 \\
media & naive & 1 & 1 & 6.161242 \\
media & naive & 2 & 1 & 1.282943 \\
media & naive & 4 & 1 & 1.541361 \\
media & secuencial & 1 & 1 & 6.036197 \\
naive\_mejor & aleatorio & 1 & 1 & 4.486226 \\
naive\_mejor & aleatorio & 2 & 1 & 2.146323 \\
naive\_mejor & aleatorio & 4 & 1 & 1.271098 \\
naive\_mejor & ciclico & 1 & 1 & 2.427398 \\
naive\_mejor & ciclico & 2 & 1 & 1.287073 \\
naive\_mejor & ciclico & 4 & 1 & 0.896508 \\
naive\_mejor & maestro\_trabajador & 1 & 1 & 2.448469 \\
naive\_mejor & maestro\_trabajador & 2 & 1 & 2.766177 \\
naive\_mejor & maestro\_trabajador & 4 & 1 & 1.499389 \\
naive\_mejor & naive & 1 & 1 & 2.431976 \\
naive\_mejor & naive & 2 & 1 & 2.546940 \\
naive\_mejor & naive & 4 & 1 & 0.043471 \\
naive\_mejor & secuencial & 1 & 1 & 2.444604 \\
naive\_peor & aleatorio & 1 & 1 & 5.998752 \\
naive\_peor & aleatorio & 2 & 1 & 3.205259 \\
naive\_peor & aleatorio & 4 & 1 & 2.074790 \\
naive\_peor & ciclico & 1 & 1 & 2.503140 \\
naive\_peor & ciclico & 2 & 1 & 1.258484 \\
naive\_peor & ciclico & 4 & 1 & 0.799426 \\
naive\_peor & maestro\_trabajador & 1 & 1 & 2.524083 \\
naive\_peor & maestro\_trabajador & 2 & 1 & 2.445595 \\
naive\_peor & maestro\_trabajador & 4 & 1 & 1.576912 \\
naive\_peor & naive & 1 & 1 & 2.467739 \\
naive\_peor & naive & 2 & 1 & 2.509108 \\
naive\_peor & naive & 4 & 1 & 3.001856 \\
naive\_peor & secuencial & 1 & 1 & 2.419464 \\
\bottomrule
\end{longtable}
}{%
  \textit{[Tabla pendiente: ejecute \texttt{scripts/run\_bench.sh} y \texttt{scripts/plot.py}.]}}

\paragraph{Capturas de evidencia.} Las salidas de consola de las corridas se conservan en
\texttt{resultados/} como evidencia de ejecución.

% ====================== REFERENCIAS ======================
\newpage
\bibliographystyle{plain}
\bibliography{refs}

\end{document}
